\chapter[Band 37: Endliche Halbgruppen]{Band 37 auf einen Blick}
\noindent\textbf{Endliche Halbgruppen}\par\medskip
Vorausgesetzt werden die Halbgruppen aus Band~28, Endlichkeit und
endliche Kardinalitäten. Die Rekonstruktion aus der großen Potenzhalbgruppe
wird nur in den ausdrücklich angegebenen Fällen behauptet.
\(A\cong B\) bezeichnet hier einen Halbgruppenisomorphismus.

\begin{BandUebersicht}
\textbf{Endlicher Grundträger}
\UebFormel{\begin{aligned}
&\FiniteSemigroupAlg A\star\\
&\quad\leftrightarrow\SemigroupAlg A\star
 \land\Finite A,\\
&X\star_{\mathcal P}Y
 =\{x\star y\mid x\in X,\ y\in Y\}.
\end{aligned}}
\UebQuellen{\FormulaRefAuto{FiniteSemigroupDef}[def];
\FormulaRefAuto{PowerSemigroupProductDef}[def]}
&
\textbf{Die große Potenzhalbgruppe}
\UebFormel{\begin{aligned}
&\FiniteSemigroupAlg A\star\ \vdash\\
&\FiniteSemigroupAlg{\PowNE A}{\star_{\mathcal P}},\\
&\FiniteCard(\PowNE A)=2^{\FiniteCard(A)}-1.
\end{aligned}}
\UebQuellen{\FormulaRefAuto{FinitePowerSemigroup}[thm];
\FormulaRefAuto{FiniteSemigroupNonemptyPowerSetCardinality}[thm]} \\
\addlinespace[.8ex]

\textbf{Index und Periode}
Für \(a\in A\) bedeutet
\(\operatorname{IP}_{A,\star}(a;\mu,\lambda)\), wobei \(r,s\in\mathbb N\):\UebFormel{\begin{aligned}
&\mu,\lambda\in\mathbb N\setminus\{0\},\\
&a^{\mu+\lambda}=a^\mu,\\
&1\le r<s<\mu+\lambda\ \vdash\ a^r\ne a^s.
\end{aligned}}
\UebQuellen{\FormulaRefAuto{FiniteSemigroupIndexPeriodPairDef}[def]}
&
\textbf{Eindeutigkeit und Potenzfolge}
In jeder endlichen Halbgruppe besitzt jedes \(a\in A\) genau ein
solches Paar. Für dieses gilt:\UebFormel{\begin{aligned}
&r\in\mathbb N,\ \mu\le r\ \vdash\\
&\qquad a^{r+\lambda}=a^r,\\
&\FiniteCard(\langle a\rangle)=\mu+\lambda-1.
\end{aligned}}
\UebQuellen{\FormulaRefAuto{FiniteSemigroupIndexPeriod}[thm];
\FormulaRefAuto{FiniteSemigroupPowerEventuallyPeriodic}[thm];
\FormulaRefAuto{FiniteSemigroupMonogenicSize}[thm]} \\
\addlinespace[.8ex]

\textbf{Monogene Halbgruppe und Idempotente}
Für \(\SemigroupAlg A\star\) und \(a\in A\):
\UebFormel{\begin{aligned}
&\langle a\rangle=\{a^n\mid
 n\in\mathbb N,\ n\ne0\},\\
&E_{A,\star}=\{e\in A\mid e\star e=e\}.
\end{aligned}}
\UebQuellen{\FormulaRefAuto{SemigroupMonogenicSubsemigroupDef}[def];
\FormulaRefAuto{SemigroupMonogenicPowerDescription}[thm];
\FormulaRefAuto{SemigroupIdempotentSetDef}[def]}
&
\textbf{Idempotente Potenzen}
\UebFormel{\begin{aligned}
&\FiniteSemigroupAlg A\star,\ a\in A\ \vdash\\
&\quad\exists n\in\mathbb N\setminus\{0\}\ 
 a^n\in E_{A,\star},\\
&\FiniteSemigroupAlg A\star,\ A\ne\varnothing\\
&\quad\vdash E_{A,\star}\ne\varnothing.
\end{aligned}}
\UebQuellen{\FormulaRefAuto{FiniteSemigroupIdempotentPower}[thm];
\FormulaRefAuto{FiniteSemigroupIdempotentExists}[thm]} \\
\addlinespace[.8ex]

\textbf{Kleinstes Ideal}
Für eine nichtleere endliche Halbgruppe \((A,\star)\) bezeichne
\(\operatorname{Id}(I)\) ein nichtleeres zweiseitiges Ideal:\UebFormel{\begin{aligned}
&\mathcal I=\{I\subseteq A\mid\operatorname{Id}(I)\},\\
&K=\iota I\bigl(\operatorname{Id}(I)\land\\
&\qquad\forall J\in\mathcal I\ (I\subseteq J)\bigr).
\end{aligned}}
\UebQuellen{\FormulaRefAuto{FiniteSemigroupIdealFamilyDef}[def];
\FormulaRefAuto{FiniteSemigroupLeastIdealDef}[def]}
&
\textbf{Existenz und Eindeutigkeit}
\UebFormel{\begin{aligned}
&\FiniteSemigroupAlg A\star,\ A\ne\varnothing\\
&\vdash\operatorname{Id}(K)\ \land\
 \forall I\in\mathcal I\ (K\subseteq I).
\end{aligned}}Damit ist \(K\) das eindeutig bestimmte kleinste nichtleere
zweiseitige Ideal.
\UebQuellen{\FormulaRefAuto{FiniteSemigroupLeastIdeal}[thm];
\FormulaRefAuto{FiniteSemigroupLeastIdealCharacterization}[thm]} \\
\addlinespace[.8ex]

\textbf{Beschriftete Multiplikationstafeln}
Für \(n=\FiniteCard(A)\) und \(e:\NatSeg n\bij A\) wird die Operation
auf \(\NatSeg n\) transportiert:\UebFormel{\begin{aligned}
&m_e(i,j)=e^{-1}(e(i)\star e(j)).
\end{aligned}}Der kanonische Code \(\operatorname{can}(A,\star)\) ist das Paar
aus \(n\) und der Menge aller zeilenweise gelesenen Tafeltupel
\(w_n(m_e)\) für sämtliche solchen Beschriftungen \(e\).
\UebQuellen{\FormulaRefAuto{FiniteSemigroupMultiplicationTableDef}[def];
\FormulaRefAuto{FiniteSemigroupTableTupleDef}[def];
\FormulaRefAuto{FiniteSemigroupCanonicalTableCodeDef}[def]}
&
\textbf{Klassifikation durch den Code}
Für endliche Halbgruppen \((A,\star)\), \((B,\diamond)\):\UebFormel{\begin{aligned}
&(A,\star)\cong(B,\diamond)\\
&\quad\leftrightarrow\
 \operatorname{can}(A,\star)
 =\operatorname{can}(B,\diamond).
\end{aligned}}Umbenennungsorbits der assoziativen Tafeln sind genau
die Isomorphieklassen.
\UebQuellen{\FormulaRefAuto{FiniteSemigroupCanonicalTableCodeComplete}[thm];
\FormulaRefAuto{FiniteSemigroupRelabelOrbitsAreIsomorphismClasses}[thm]} \\
\addlinespace[.8ex]

\textbf{Einermengenschicht}
\UebFormel{\begin{aligned}
&\operatorname{Sing}(A)=\{\{a\}\mid a\in A\}.
\end{aligned}}Sei \(F\) ein Isomorphismus der großen Potenzhalbgruppen von
\((A,\star)\) und \((B,\diamond)\).
\UebQuellen{\FormulaRefAuto{SemigroupSingletonCopyDef}[def];
\FormulaRefAuto{SemigroupIsoDef}[def]}
&
\textbf{Ein rekonstruierbarer Teilfall}
\UebFormel{\begin{aligned}
&F[\operatorname{Sing}(A)]
 =\operatorname{Sing}(B)\\
&\qquad\vdash (A,\star)\cong(B,\diamond).
\end{aligned}}Für endliche Halbgruppen mit
\(\FiniteCard(A)=\FiniteCard(B)=2\) genügt bereits die Existenz von \(F\).
\UebQuellen{\FormulaRefAuto{PowerSemigroupSingletonLayerReconstruction}[thm];
\FormulaRefAuto{FinitePowerSemigroupReconstructionOrderTwo}[thm]} \\
\addlinespace[.8ex]

\textbf{Endliche Suche}
Auf dem festen Träger \(\NatSeg n\) werden Operationen
\(m:(\NatSeg n)^2\to\NatSeg n\) mit\UebFormel{\begin{aligned}
&m(m(i,j),k)=m(i,m(j,k))\\
&\hfill(i,j,k\in\NatSeg n)
\end{aligned}}betrachtet und anschließend nach Umbenennung zusammengefasst.
\UebQuellen{\FormulaRefAuto{FiniteSemigroupAssociativeOperationsDef}[def];
\FormulaRefAuto{FiniteSemigroupRelabelledOperationDef}[def]}
&
\textbf{Reichweite des Verfahrens}
Die Tafelkodierung erlaubt die vollständige Prüfung einer fest
vorgegebenen endlichen Ordnung. Sie beweist nicht die allgemeine
Implikation\UebFormel{\begin{aligned}
&(\PowNE A,\star_{\mathcal P})
 \cong(\PowNE B,\diamond_{\mathcal P})\\
&\qquad\Longrightarrow (A,\star)\cong(B,\diamond).
\end{aligned}}Diese bleibt im allgemeinen endlichen Fall eine Vermutung.
\UebQuellen{\FormulaRefAuto{FiniteSemigroupCanonicalTableCodeComplete}[thm];
\FormulaRefAuto{FinitePowerSemigroupReconstructionOrderTwo}[thm]} \\
\addlinespace[.8ex]
\end{BandUebersicht}

